\textbf{Phần 1. Trắc nghiệm lựa chọn}
\setcounterpageref{ex}{0}

\begin{ex} %[5D2N12]
Viết số đo $2\text{ m } 15\text{ cm}$ dưới dạng số thập phân có đơn vị đo là mét ta được
\choice
{$21{,}5\text{ m}$}
{$2{,}015\text{ m}$}
{\True $2{,}15\text{ m}$}
{$0{,}215\text{ m}$}
\loigiai{
Ta có: $2\text{ m } 15\text{ cm} = 2\dfrac{15}{100}\text{ m} = 2{,}15\text{ m}$.\\
\textbf{Đáp án của câu hỏi là C}
}
\end{ex}

\begin{ex} %[5D2N12]
Viết số đo $3\text{ m } 8\text{ cm}$ dưới dạng số thập phân có đơn vị mét ta được
\choice
{$3{,}8\text{ m}$}
{\True $3{,}08\text{ m}$}
{$3{,}80\text{ m}$}
{$0{,}38\text{ m}$}
\loigiai{
Ta có: $3\text{ m } 8\text{ cm} = 3\dfrac{8}{100}\text{ m} = 3{,}08\text{ m}$.\\
\textbf{Đáp án của câu hỏi là B}
}
\end{ex}

\begin{ex} %[5D2H12]
Viết số đo $7\text{ km } 80\text{ m}$ dưới dạng số thập phân có đơn vị ki-lô-mét ta được
\choice
{$7{,}8\text{ km}$}
{$7{,}008\text{ km}$}
{$78{,}0\text{ km}$}
{\True $7{,}08\text{ km}$}
\loigiai{
Ta có: $7\text{ km } 80\text{ m} = 7\dfrac{80}{1\,000}\text{ km} = 7{,}08\text{ km}$.\\
\textbf{Đáp án của câu hỏi là D}
}
\end{ex}

\begin{ex} %[5D2H12]
Viết số đo $8\text{ kg } 75\text{ g}$ dưới dạng số thập phân có đơn vị ki-lô-gam ta được
\choice
{$8{,}75\text{ kg}$}
{\True $8{,}075\text{ kg}$}
{$8{,}750\text{ kg}$}
{$80{,}75\text{ kg}$}
\loigiai{
Ta có: $8\text{ kg } 75\text{ g} = 8\dfrac{75}{1\,000}\text{ kg} = 8{,}075\text{ kg}$.\\
\textbf{Đáp án của câu hỏi là B}
}
\end{ex}

\begin{ex} %[5D2H12]
Viết số đo $2\text{ tấn } 325\text{ kg}$ dưới dạng số thập phân với đơn vị tấn ta được
\choice
{\True $2{,}325\text{ tấn}$}
{$23{,}25\text{ tấn}$}
{$2{,}0325\text{ tấn}$}
{$2{,}3250\text{ tấn}$}
\loigiai{
Ta có: $2\text{ tấn } 325\text{ kg} = 2\dfrac{325}{1\,000}\text{ tấn} = 2{,}325\text{ tấn}$.\\
\textbf{Đáp án của câu hỏi là A}
}
\end{ex}

\begin{ex} %[5D2H12]
Viết số đo diện tích $1\text{ m}^2\ 60\text{ dm}^2$ dưới dạng số thập phân có đơn vị mét vuông ta được
\choice
{$1{,}06\text{ m}^2$}
{\True $1{,}6\text{ m}^2$}
{$16{,}0\text{ m}^2$}
{$1{,}006\text{ m}^2$}
\loigiai{
Ta có: $1\text{ m}^2\ 60\text{ dm}^2 = 1\dfrac{60}{100}\text{ m}^2 = 1{,}60\text{ m}^2 = 1{,}6\text{ m}^2$.\\
\textbf{Đáp án của câu hỏi là B}
}
\end{ex}

\begin{ex} %[5D2H12]
Viết số đo diện tích $3\text{ m}^2\ 6\text{ dm}^2$ dưới dạng số thập phân có đơn vị mét vuông ta được
\choice
{$3{,}6\text{ m}^2$}
{$3{,}006\text{ m}^2$}
{\True $3{,}06\text{ m}^2$}
{$36{,}0\text{ m}^2$}
\loigiai{
Ta có: $3\text{ m}^2\ 6\text{ dm}^2 = 3\dfrac{6}{100}\text{ m}^2 = 3{,}06\text{ m}^2$.\\
\textbf{Đáp án của câu hỏi là C}
}
\end{ex}

\begin{ex} %[5D2H12]
Viết số đo diện tích $120\text{ dm}^2$ dưới dạng số thập phân có đơn vị mét vuông ta được
\choice
{$12{,}0\text{ m}^2$}
{$0{,}12\text{ m}^2$}
{$1{,}02\text{ m}^2$}
{\True $1{,}2\text{ m}^2$}
\loigiai{
Ta có: $120\text{ dm}^2 = \dfrac{120}{100}\text{ m}^2 = 1\dfrac{20}{100}\text{ m}^2 = 1{,}2\text{ m}^2$.\\
\textbf{Đáp án của câu hỏi là D}
}
\end{ex}

\begin{ex} %[5D2H12]
Viết số đo dung tích $5\text{ l } 75\text{ ml}$ dưới dạng số thập phân có đơn vị lít ta được
\choice
{\True $5{,}075\text{ l}$}
{$5{,}75\text{ l}$}
{$5{,}750\text{ l}$}
{$50{,}75\text{ l}$}
\loigiai{
Ta có: $5\text{ l } 75\text{ ml} = 5\dfrac{75}{1\,000}\text{ l} = 5{,}075\text{ l}$.\\
\textbf{Đáp án của câu hỏi là A}
}
\end{ex}

\begin{ex} %[5D2H12]
Một con thỏ cân nặng $6\text{ kg } 75\text{ g}$, một con ngỗng cân nặng $6\,100\text{ g}$, một con mèo cân nặng $6{,}095\text{ kg}$. Con vật nào nặng nhất?
\choice
{Thỏ}
{Mèo}
{Ba con nặng bằng nhau}
{\True Ngỗng}
\loigiai{
Đổi về cùng đơn vị ki-lô-gam:
\begin{itemize}
    \item Thỏ: $6\text{ kg } 75\text{ g} = 6{,}075\text{ kg}$.
    \item Ngỗng: $6\,100\text{ g} = 6{,}1\text{ kg} = 6{,}100\text{ kg}$.
    \item Mèo: $6{,}095\text{ kg}$.
\end{itemize}
Vì $6{,}075 < 6{,}095 < 6{,}100$ nên con ngỗng nặng nhất.\\
\textbf{Đáp án của câu hỏi là D}
}
\end{ex}

\begin{ex} %[5D2V12]
Có ba bức tranh tường: Tranh về bảo vệ môi trường có diện tích $5{,}3\text{ m}^2$; Tranh về an toàn giao thông có diện tích $5\text{ m}^2\ 8\text{ dm}^2$; Tranh về phòng chống Covid có diện tích $5\text{ m}^2\ 9\text{ dm}^2$. Bức tranh có diện tích bé nhất là
\choice
{Tranh về bảo vệ môi trường}
{Tranh về phòng chống Covid}
{\True Tranh về an toàn giao thông}
{Ba bức tranh có diện tích bằng nhau}
\loigiai{
Đổi về cùng đơn vị mét vuông:
\begin{itemize}
    \item Bảo vệ môi trường: $5{,}3\text{ m}^2 = 5{,}30\text{ m}^2$.
    \item An toàn giao thông: $5\text{ m}^2\ 8\text{ dm}^2 = 5{,}08\text{ m}^2$.
    \item Phòng chống Covid: $5\text{ m}^2\ 9\text{ dm}^2 = 5{,}09\text{ m}^2$.
\end{itemize}
Vì $5{,}08 < 5{,}09 < 5{,}30$ nên bức tranh về an toàn giao thông có diện tích bé nhất.\\
\textbf{Đáp án của câu hỏi là C}
}
\end{ex}

\begin{ex} %[5D2V12]
Đoạn đường $AB$ dài $1{,}2\text{ km}$, đoạn đường $AC$ dài $1\text{ km } 75\text{ m}$. Hỏi đoạn đường nào dài hơn và dài hơn bao nhiêu mét?
\choice
{\True $AB$ dài hơn $AC$ là $125\text{ m}$}
{$AC$ dài hơn $AB$ là $125\text{ m}$}
{$AB$ dài hơn $AC$ là $25\text{ m}$}
{$AC$ dài hơn $AB$ là $75\text{ m}$}
\loigiai{
Đổi ra mét: $AB = 1{,}2\text{ km} = 1\,200\text{ m}$; $AC = 1\text{ km } 75\text{ m} = 1\,075\text{ m}$.\\
Đoạn đường $AB$ dài hơn đoạn đường $AC$ là: $1\,200 - 1\,075 = 125\text{ (m)}$.\\
\textbf{Đáp án của câu hỏi là A}
}
\end{ex}

\textbf{Phần 2. Trắc nghiệm đúng sai}
\setcounterpageref{ex}{0}

\begin{ex} %[5D2H12]
Xét tính đúng sai của các phép đổi đơn vị đo sau:
\choiceTF
{\True $4\text{ km } 500\text{ m} = 4{,}5\text{ km}$.}
{$1\text{ tấn } 5\text{ tạ} = 1{,}05\text{ tấn}$.}
{\True $8\text{ m}^2\ 75\text{ dm}^2 = 8{,}75\text{ m}^2$.}
{$2\text{ dm}^2\ 5\text{ cm}^2 = 2{,}5\text{ dm}^2$.}
\loigiai{
\begin{itemize}
    \item $4\text{ km } 500\text{ m} = 4\dfrac{500}{1\,000}\text{ km} = 4{,}5\text{ km}$.
    \item $1\text{ tấn } 5\text{ tạ} = 1\dfrac{5}{10}\text{ tấn} = 1{,}5\text{ tấn} \neq 1{,}05\text{ tấn}$.
    \item $8\text{ m}^2\ 75\text{ dm}^2 = 8\dfrac{75}{100}\text{ m}^2 = 8{,}75\text{ m}^2$.
    \item $2\text{ dm}^2\ 5\text{ cm}^2 = 2\dfrac{5}{100}\text{ dm}^2 = 2{,}05\text{ dm}^2 \neq 2{,}5\text{ dm}^2$.
\end{itemize}
\begin{itemize}
    \item a) \textbf{Đúng}.
    \item b) \textbf{Sai}: $1\text{ tấn } 5\text{ tạ} = 1{,}5\text{ tấn}$.
    \item c) \textbf{Đúng}.
    \item d) \textbf{Sai}: $2\text{ dm}^2\ 5\text{ cm}^2 = 2{,}05\text{ dm}^2$.
\end{itemize}
}
\end{ex}

\begin{ex} %[5D2V12]
Hình $A$ có diện tích là $4\text{ cm}^2\ 15\text{ mm}^2$ và Hình $B$ có diện tích là $3{,}95\text{ cm}^2$. Xét tính đúng sai của các khẳng định sau:
\choiceTF
{\True Đổi diện tích hình $A$ sang xăng-ti-mét vuông ta được $4{,}15\text{ cm}^2$.}
{Hình $B$ có diện tích lớn hơn hình $A$.}
{\True Diện tích hình $A$ lớn hơn diện tích hình $B$ là $0{,}2\text{ cm}^2$.}
{Tổng diện tích của cả hai hình là $8\text{ cm}^2$.}
\loigiai{
\begin{itemize}
    \item Hình $A$: $4\text{ cm}^2\ 15\text{ mm}^2 = 4\dfrac{15}{100}\text{ cm}^2 = 4{,}15\text{ cm}^2$.
    \item So sánh: $4{,}15\text{ cm}^2 > 3{,}95\text{ cm}^2$ nên hình $A$ có diện tích lớn hơn hình $B$.
    \item Hiệu diện tích: $4{,}15 - 3{,}95 = 0{,}2\text{ cm}^2$.
    \item Tổng diện tích: $4{,}15 + 3{,}95 = 8{,}1\text{ cm}^2 \neq 8\text{ cm}^2$.
\end{itemize}
\begin{itemize}
    \item a) \textbf{Đúng}.
    \item b) \textbf{Sai}: Hình $A$ lớn hơn hình $B$.
    \item c) \textbf{Đúng}.
    \item d) \textbf{Sai}: Tổng diện tích là $8{,}1\text{ cm}^2$.
\end{itemize}
}
\end{ex}

\textbf{Phần 3. Trả lời ngắn}
\setcounterpageref{ex}{0}

\begin{ex} %[5D2H12]
Viết số đo $456\text{ m}$ dưới dạng số thập phân có đơn vị ki-lô-mét. Kết quả là bao nhiêu? (Nhập dưới dạng số thập phân).
\shortans{0,456}
\loigiai{
\textbf{Đáp số:} $0{,}456$\\
$456\text{ m} = \dfrac{456}{1\,000}\text{ km} = 0{,}456\text{ km}$.
}
\end{ex}

\begin{ex} %[5D2H12]
Viết số đo $1\,450\text{ kg}$ dưới dạng số thập phân có đơn vị tấn (viết dưới dạng gọn nhất).
\shortans{1,45}
\loigiai{
\textbf{Đáp số:} $1{,}45$\\
$1\,450\text{ kg} = \dfrac{1\,450}{1\,000}\text{ tấn} = 1{,}45\text{ tấn}$.
}
\end{ex}

\begin{ex} %[5D2H12]
Viết số đo $85\text{ cm}^2$ dưới dạng số thập phân có đơn vị đề-xi-mét vuông.
\shortans{0,85}
\loigiai{
\textbf{Đáp số:} $0{,}85$\\
$85\text{ cm}^2 = \dfrac{85}{100}\text{ dm}^2 = 0{,}85\text{ dm}^2$.
}
\end{ex}

\begin{ex} %[5D2H12]
Viết số đo $3\,452\text{ ml}$ dưới dạng số thập phân có đơn vị lít.
\shortans{3,452}
\loigiai{
\textbf{Đáp số:} $3{,}452$\\
$3\,452\text{ ml} = \dfrac{3\,452}{1\,000}\text{ l} = 3{,}452\text{ l}$.
}
\end{ex}

\begin{ex} %[5D2V12]
Một hộp sữa và $1$ quả cân $250\text{ g}$ ở một đĩa cân thăng bằng với $1$ quả cân $1\text{ kg}$ và $1$ quả cân $500\text{ g}$ ở đĩa cân bên kia. Hỏi hộp sữa đó cân nặng bao nhiêu ki-lô-gam? (Nhập dưới dạng số thập phân).
\shortans{1,25}
\loigiai{
\textbf{Đáp số:} $1{,}25$\\
Khối lượng hộp sữa là: $(1\,000 + 500) - 250 = 1\,250\text{ (g)} = 1{,}25\text{ (kg)}$.
}
\end{ex}

\begin{ex} %[5D2V12]
Con ốc sên bò theo đường gấp khúc $ABC$, trong đó đoạn $AB = 1{,}8\text{ m}$ và đoạn $BC = 2\text{ m } 15\text{ cm}$. Hỏi con ốc sên đã bò tất cả bao nhiêu mét? (Nhập dưới dạng số thập phân).
\shortans{3,95}
\loigiai{
\textbf{Đáp số:} $3{,}95$\\
Đổi $2\text{ m } 15\text{ cm} = 2{,}15\text{ m}$.\\
Tổng quãng đường ốc sên bò là: $1{,}8 + 2{,}15 = 3{,}95\text{ (m)}$.
}
\end{ex}
